\chapter{Quantum Computation and the Limits of Randomness}
\label{ch:quantum}

Randomness is one resource that might extend polynomial time computation; quantum mechanics is another. The class $\mathrm{BQP}$ of problems solvable in polynomial time by a quantum computer with bounded error contains $\BPP$ and, by Shor\index{Shor's algorithm}'s algorithm, factoring and the discrete logarithm, which are not known to be in $\BPP$. This chapter describes the class, its relation to $\NP$, and what quantum methods do and do not offer for counting, because the Monte Carlo program of Part IV has a quantum variant and it is worth knowing its limits.

\section{The class BQP}

A quantum computer acts on a register of $n$ qubits, whose state is a unit vector in $\mathbb{C}^{2^n}$, by a sequence of unitary operations each acting on a constant number of qubits, followed by a measurement that returns a basis state with probability equal to the squared modulus of its amplitude. A language is in $\mathrm{BQP}$ if there is a uniformly generated family of polynomial size quantum circuits that accepts each string in the language with probability at least $\tfrac23$ and each string outside with probability at most $\tfrac13$. The model was formalised by Bernstein and Vazirani in 1993 and Yao in 1993, and the inclusions
\[
\BPP\index{BPP} \subseteq \mathrm{BQP} \subseteq \mathrm{PP} \subseteq \mathrm{PSPACE}
\]
are known. The first holds because a quantum circuit can simulate a coin. The last two hold because the acceptance probability of a quantum circuit is a sum of exponentially many terms, each computable in polynomial time, and such a sum can be compared to a threshold in $\mathrm{PP}$ by Adleman, DeMarrais and Huang and in $\mathrm{PSPACE}$ by direct simulation.

\section{Shor and Grover}

Shor's algorithm of 1994 factors an $n$ bit integer in time $O(n^3)$ on a quantum computer, by reducing factoring to period finding and finding the period with the quantum Fourier transform. It places factoring in $\mathrm{BQP}$ and is the strongest evidence that $\mathrm{BQP} \neq \BPP$, since no classical algorithm for factoring is known to run faster than $2^{O(n^{1/3} \log^{2/3} n)}$. It also breaks RSA and the Diffie Hellman and elliptic curve systems, which is why post quantum cryptography exists.

Grover\index{Grover's algorithm}'s algorithm of 1996 searches an unstructured list of $M$ items for a marked one in $O(\sqrt M)$ queries, against $M$ classically, and Bennett, Bernstein, Brassard and Vazirani showed that $\sqrt M$ is optimal for black box search. The consequence for $\NP$ is that a quantum computer can solve $\mathrm{SAT}$ on $n$ variables in time $2^{n/2} \mathrm{poly}(n)$ by searching the assignments, a quadratic speedup over brute force, and that no black box quantum algorithm does better. This is the main evidence that $\NP \not\subseteq \mathrm{BQP}$: relative to a random oracle, $\NP$ is not in $\mathrm{BQP}$, and the structure of $\NP$ complete problems does not seem to offer the periodicity that Shor exploits.

\section{Quantum counting and amplitude estimation}

For the counting problems of this book the relevant quantum primitive is amplitude estimation\index{amplitude estimation}, due to Brassard, H{\o}yer, Mosca and Tapp in 2000. Given a procedure that prepares a state with amplitude $\sqrt{p}$ on a marked subspace, amplitude estimation returns $p$ to within $\epsilon$ using $O(1/\epsilon)$ applications of the procedure, against $O(1/\epsilon^2)$ samples for classical Monte Carlo. This is a quadratic speedup in the accuracy, and it applies directly to the rare event estimation of Chapter~\ref{ch:rare}: the probability $\rho = \PP(S = 0)$ under the uniform distribution can be estimated to relative error $\epsilon$ with $O(1 / (\epsilon \sqrt\rho))$ queries to the energy function, against $O(1 / (\epsilon^2 \rho))$ classically. Montanaro's 2015 paper gives the general Monte Carlo speedup, and Szegedy's 2004 quantum walk gives a quadratic speedup in the hitting time of a reversible Markov chain. A square root speedup for sampling is known for annealing sequences of chains whose consecutive stationary distributions overlap substantially (Somma, Boixo, Barnum and Knill 2008; Wocjan and Abeyesinghe 2008), which is the setting of the ladder; we do not claim it for the mixing time of a single chain, and below the square root speedup for the conditioned draws is conditional on that overlap hypothesis.

Put together, the quantum version of the vertical walk estimates each rung of the ladder with quadratically fewer samples and quadratically faster conditioned draws (under the overlap hypothesis above). The ladder itself, its length and the number of rungs, is unchanged, since it is a property of the staircase and not of the sampler. For the queens the classical cost is $e^{0.942 N}$ for naive Monte Carlo and polynomial in $N$ times the mixing time\index{mixing time} for the walk; the quantum cost is $e^{0.471 N}$ for naive amplitude estimation and polynomial times the square root of the mixing time for the walk, under the same hypothesis. If the mixing time is exponential, so is the quantum walk, with at best half the exponent under that hypothesis. A quadratic speedup does not turn an exponential algorithm into a polynomial one, and $\sharpP$ complete problems are not believed to be in $\mathrm{BQP}$ even approximately, since approximate counting of $\NP$ complete problems in $\mathrm{BQP}$ would place $\NP$ in $\mathrm{BQP}$.

\section{Quantum Bayesian computation}

Polson, Sokolov and Xu\index{Polson, Sokolov and Xu} set out in 2023 the quantum version of the Bayesian computation that this book turns into counting. Their observation is that the primitives of Bayesian inference, sampling from a posterior, computing a normalising constant, and estimating a posterior expectation, are each a sum over an exponentially large space of a quantity computable term by term, and that the quantum algorithms of Grover and of Brassard, H{\o}yer, Mosca and Tapp apply to such sums with a quadratic gain. Specifically, they show that the evidence $Z = \EE_P[L(X)]$ under the prior $P$ can be estimated to relative error $\epsilon$ with $O(1/\epsilon)$ quantum queries to an oracle that prepares the prior and evaluates the likelihood, against $O(1/\epsilon^2)$ classical samples; that posterior sampling by rejection from the prior, with acceptance probability $p$, costs $O(1/\sqrt p)$ quantum oracle calls by amplitude amplification against $O(1/p)$ classically; and that these gains compose with the quantile reordering of vertical likelihood, so that each rung of a ladder is estimated with the square root of the classical cost. They also give the quantum form of the von Neumann rejection sampler of Chapter~\ref{ch:shannon}, in which the accept reject step becomes a reflection and the expected number of candidates per accepted sample drops from $c$ to $O(\sqrt c)$.

For the queens the arithmetic is as follows. Naive rejection from uniform permutations has acceptance probability $\rho_N \approx e^{-0.942 N}$ and classical cost $e^{0.942 N}$ per solution; the quantum rejection sampler of Polson, Sokolov and Xu costs $e^{0.471 N}$. The classical ladder with a rapidly mixing level set chain costs about $N^{4.8} / \epsilon^2$ by the measured exponent of Chapter~\ref{ch:experiments2}; the quantum ladder, with quantum walks for the conditioned draws (under the overlap hypothesis) and amplitude estimation for the ratios, would cost, if the square root speedups combine as assumed, about $N^{3.4} / \epsilon$ by halving the autocorrelation time exponent and the accuracy exponent. On conditioned boards where the chain is torpid with relaxation time $e^{cN}$, the quantum ladder would have relaxation cost $e^{cN/2}$ under that hypothesis, still exponential. The paper's conclusion, that quantum Bayesian computation gives a quadratic and never an exponential speedup on the generic sum, is the conclusion of this chapter for counting: the staircase is unchanged, every rung is cheaper by a square root, and the question whether the bottom rung can be reached in polynomial time is the same question.

What the quantum Bayesian view adds to the program is a precise statement of what a quantum computer would and would not change in the main equivalence of Chapter~\ref{ch:theorems}. It would not change the equivalence of $\NP = \RP$ with the polynomial samplability of the queens level sets, since that is a statement about classical randomized computation. One would expect a parallel equivalence, which we state as a conjecture and do not prove: $\NP \subseteq \mathrm{BQP}$ if and only if the queens level sets admit a polynomial time quantum sampler (a quantum algorithm that outputs a state within inverse polynomial trace distance of the uniform superposition over the level set), by the same chain of reductions with Bellare, Goldreich and Petrank's uniform generation run with a $\mathrm{BQP}$ oracle in place of an $\NP$ oracle. Making this precise needs $\mathrm{BQP}^{\mathrm{BQP}} = \mathrm{BQP}$ and an error analysis for a bounded error oracle inside the generator, with a strict time bound and a failure symbol in place of an expected time bound. And it would locate, in the quantum world, the same obstruction: Grover's lower bound says that a rung of density $\rho$ cannot be hit with fewer than order $1/\sqrt{\rho}$ queries by a black box quantum search, so the obstruction is at the sparse bottom rung, and the oracle separation of $\NP$ from $\mathrm{BQP}$ is consistent with there being no structure in the queens that a quantum walk could exploit beyond the square root.

\section{Randomness in the quantum world}

Is $\mathrm{BQP} = \mathrm{P}$? Nobody believes it, because of Shor. Is $\mathrm{BQP}$ derandomizable in the sense that $\BPP$ is believed to be? Not in the same way: the randomness in a quantum algorithm is in the measurement, not in a coin that could be replaced by a pseudorandom generator\index{pseudorandom generator}, and no derandomization theorem is known for $\mathrm{BQP}$. The hardness versus randomness paradigm applies to classical randomness and leaves the quantum case untouched. Aaronson has argued that the relation between $\mathrm{BQP}$ and the polynomial hierarchy is the natural quantum analogue of $\Pc$ versus $\NP$, and Raz and Tal showed in 2018 that relative to an oracle $\mathrm{BQP}$ is not in $\PH$, so that quantum computation is not merely randomness with better constants.

\section{For the program}

The Monte Carlo program of Part IV is classical. Its quantum variant replaces each conditioned draw by a quantum walk and each ratio estimate by amplitude estimation, and gains a square root in each, the first under the overlap hypothesis stated above. The staircase is the same, the hardness is in the same place, and the conclusion is the same: the walk is polynomial if and only if the level sets can be traversed, and for an $\NP$ complete energy they cannot be, in the quantum world any more than the classical, unless $\NP \subseteq \mathrm{BQP}$, which Grover's lower bound and the oracle results make unlikely. Quantum mechanics at best halves the exponent, under the overlap hypothesis, and leaves the question alone.
